\setcounter{section}{6}

  \zwischenueberschrift{Übungsaufgaben}

\inputaufgabe
{}
{

Es sei $Y$ der
\definitionsverweis {Graph}{}{}
zur Funktion

\mavergleichskette
{\vergleichskette
{ f(x,y)
}
{ = }{ x^2-y^2
}
{  }{
}
{    }{
}
{   }{
}
}
{}{}{.}
\aufzaehlungdrei{Zeige, dass durch

\mavergleichskettedisp
{\vergleichskette
{ F(x,y,z)
}
{ =} { \begin{pmatrix}  y  \\ x \\  0   \end{pmatrix}
}
{ } {
}
{ } {
}
{ } {
}
}
{}{}{}
ein
\definitionsverweis {tangentiales Vektorfeld}{}{}
auf $Y$ gegeben ist.
}{Es sei

\mavergleichskette
{\vergleichskette
{ P
}
{ = }{ \begin{pmatrix}  2  \\ 3 \\  -5   \end{pmatrix}
}
{ \in }{ Y
}
{    }{
}
{   }{
}
}
{}{}{.}
Berechne
\mathl{\nabla_v F}{} für

\mavergleichskette
{\vergleichskette
{ v
}
{ = }{ \begin{pmatrix}  2  \\ 1 \\  2   \end{pmatrix}
}
{ = }{ T_PY
}
{    }{
}
{   }{
}
}
{}{}{.}
}{Überprüfe, dass $\nabla_v F$ tangential an $Y$ in $P$ ist.
}

}
{} {}

\inputaufgabegibtloesung
{}
{

Wir betrachten auf der zweidimensionalen
\definitionsverweis {Einheitssphäre}{}{}
$Y$ das
\definitionsverweis {tangentiale Vektorfeld}{}{}

\mavergleichskettedisp
{\vergleichskette
{ F \begin{pmatrix}  x  \\ y \\ z   \end{pmatrix}
}
{ =} { \begin{pmatrix}  y  \\ -x\\  0   \end{pmatrix}
}
{ } {
}
{ } {
}
{ } {
}
}
{}{}{.}
\aufzaehlungdrei{Bestimme die
\definitionsverweis {Jacobi-Matrix}{}{}
zu $F$ auf dem $\R^3$.
}{Bestimme für den Punkt

\mavergleichskette
{\vergleichskette
{ P
}
{ = }{ \begin{pmatrix}  0  \\ 1 \\  0   \end{pmatrix}
}
{ \in }{ Y
}
{    }{
}
{   }{
}
}
{}{}{}
eine Matrix für die Abbildung
\maabbeledisp {} {T_PY} { T_PY
} { v } { \nabla_v F
} {,}
bezüglich einer geeigneten
\definitionsverweis {Basis}{}{.}
}{Bestimme für den Punkt

\mavergleichskette
{\vergleichskette
{ Q
}
{ = }{ \begin{pmatrix}  0  \\ 0\\  1   \end{pmatrix}
}
{ \in }{ Y
}
{    }{
}
{   }{
}
}
{}{}{}
eine Matrix für die Abbildung
\maabbeledisp {} {T_QY} { T_QY
} { v } { \nabla_v F
} {,}
bezüglich einer geeigneten
\definitionsverweis {Basis}{}{.}
}

}
{} {}

\inputaufgabe
{}
{

Es sei

\mavergleichskette
{\vergleichskette
{ Y
}
{ \subseteq }{ W
}
{ \subseteq }{ \R^n
}
{    }{
}
{   }{
}
}
{}{}{,}
$W$
\definitionsverweis {offen}{}{,}
eine
\definitionsverweis {differenzierbare Hyperfläche}{}{}
und sei
\maabb {\gamma} { I } { Y
} {}
eine
\definitionsverweis {differenzierbare Kurve}{}{.}
Zeige die folgenden Aussagen.
\aufzaehlungzwei {Zu einem längs $\gamma$
\definitionsverweis {parallelen Vektorfeld}{}{}
$F$ und

\mavergleichskette
{\vergleichskette
{ a
}
{ \in }{ \R
}
{  }{
}
{    }{
}
{   }{
}
}
{}{}{}
ist auch
\mathl{aF}{} ein paralleles Vektorfeld.
} {Zu längs $\gamma$ parallelen Vektorfelder
$F , G$ ist auch
\mathl{F +G}{} ein paralleles Vektorfeld.
}

}
{} {}

\inputaufgabe
{}
{

Es sei

\mavergleichskette
{\vergleichskette
{ Y
}
{ \subseteq }{ \R^n
}
{  }{
}
{    }{
}
{   }{
}
}
{}{}{}
eine
\definitionsverweis {Hyperebene}{}{.}
Zeige, dass jedes
\definitionsverweis {tangentiale Vektorfeld}{}{}
$F$
\zusatzklammer {das einem Vektorfeld auf

\mavergleichskette
{\vergleichskette
{ Y
}
{ \cong }{\R^{n-1}
}
{  }{
}
{    }{
}
{   }{
}
}
{}{}{}
entspricht} {} {}
\definitionsverweis {parallel}{}{}
\zusatzklammer {längs jeder Kurve} {} {}
ist.

}
{} {}

\inputaufgabe
{}
{

Es sei

\mavergleichskettedisp
{\vergleichskette
{ Y
}
{ =} { \left\{  \begin{pmatrix}  x  \\ y \\ z   \end{pmatrix}   \mid   x^2+y^2+z^2 = 1        \right\}
}
{ } {
}
{ } {
}
{ } {
}
}
{}{}{}
die zweidimensionale
\definitionsverweis {Einheitssphäre}{}{}
und
\maabbele {\gamma} {[0,2\pi]} {Y
} { t } { \begin{pmatrix}  \cos  t  \\ 0 \\  \sin  t   \end{pmatrix}
} {.}
Bestimme, welche der folgenden Vektorfelder
\maabbdisp {F} {[0,2\pi]} { \R^3
} {}
längs $\gamma$ tangential an $Y$  und welche
\definitionsverweis {parallel}{}{}
sind.
\aufzaehlungfuenf{
\mavergleichskette
{\vergleichskette
{ F (t)
}
{ = }{ \begin{pmatrix}  2  \\ 0 \\  1   \end{pmatrix}
}
{  }{
}
{    }{
}
{   }{
}
}
{}{}{.}
}{
\mavergleichskette
{\vergleichskette
{ F (t)
}
{ = }{ \begin{pmatrix}  0  \\ 5 \\  0   \end{pmatrix}
}
{  }{
}
{    }{
}
{   }{
}
}
{}{}{.}
}{
\mavergleichskette
{\vergleichskette
{ F (t)
}
{ = }{ \begin{pmatrix}  \cos  t  \\ 0 \\  \sin  t   \end{pmatrix}
}
{  }{
}
{    }{
}
{   }{
}
}
{}{}{.}
}{
\mavergleichskette
{\vergleichskette
{ F (t)
}
{ = }{ \begin{pmatrix}  - \sin  t  \\ 0 \\  \cos  t   \end{pmatrix}
}
{  }{
}
{    }{
}
{   }{
}
}
{}{}{.}
}{
\mavergleichskette
{\vergleichskette
{ F (t)
}
{ = }{ \begin{pmatrix}  t  \\e^t\\  \sin  t   \end{pmatrix}
}
{  }{
}
{    }{
}
{   }{
}
}
{}{}{.}
}

}
{} {}

\inputaufgabegibtloesung
{}
{

Es sei

\mavergleichskette
{\vergleichskette
{ Y
}
{ \subseteq }{ W
}
{ \subseteq }{ \R^3
}
{    }{
}
{   }{
}
}
{}{}{,}
$W$
\definitionsverweis {offen}{}{,}
eine
\definitionsverweis {differenzierbare Fläche}{}{}
versehen mit einem
\definitionsverweis {Einheitsnormalenfeld}{}{}
$N$. Es sei
\maabb {\gamma} { I } { Y
} {}
eine stetig differenzierbare Kurve und sei
\maabbdisp {F} { I } { \R^3
} {}
ein differenzierbares
\definitionsverweis {tangentiales Vektorfeld}{}{}
längs $\gamma$, das
\definitionsverweis {parallel}{}{}
sei. Zeige, dass dann auch das Vektorfeld
\maabbeledisp {} { I } { \R^3
} { t } { N(\gamma(t)) \times F(t)
} {,}
parallel ist, wobei $\times$ das
\definitionsverweis {Kreuzprodukt}{}{}
im $\R^3$ bezeichnet.

}
{} {}

\inputaufgabe
{}
{

Es sei

\mavergleichskettedisp
{\vergleichskette
{ h(x,y,z)
}
{ =} { x^3-2y^4-3z^5
}
{ } {
}
{ } {
}
{ } {
}
}
{}{}{}
und

\mavergleichskette
{\vergleichskette
{ Y
}
{ = }{ h^{-1}(1)
}
{  }{
}
{    }{
}
{   }{
}
}
{}{}{.}
Es sei
\maabbeledisp {\gamma} { \R} {Y
} { t } { \begin{pmatrix}  \sqrt[3]{ 1+3t^5  } \\ 0 \\ t   \end{pmatrix}
} {,}
ein differenzierbarer Weg auf $Y$. Bestimme die gewöhnliche Differentialgleichung in der Form von
[[Hyperfläche/Kurve/Anfangstangentenvektor/Paralleles Vektorfeld/Differentialgleichung/Explizite Gestalt/Bemerkung|Bemerkung 6.10]]
für die
\definitionsverweis {parallelen Vektorfelder}{}{}
längs $\gamma$.

}
{} {}

\inputaufgabe
{}
{

Es sei

\mavergleichskette
{\vergleichskette
{ Y
}
{ \subseteq }{ \R^n
}
{  }{
}
{    }{
}
{   }{
}
}
{}{}{}
eine
\definitionsverweis {Hyperebene}{}{.}
Zeige, dass der
\definitionsverweis {Paralleltransport}{}{}
zu jeder Kurve auf $Y$ die Identität ist, wenn man die
\definitionsverweis {Tangentialräume}{}{}
$T_PY$ mit $Y$ identifiziert.

}
{} {}

\inputaufgabegibtloesung
{}
{

Es sei

\mavergleichskette
{\vergleichskette
{ Y
}
{ \subseteq }{ W
}
{ \subseteq }{ \R^n
}
{    }{
}
{   }{
}
}
{}{}{,}
$W$
\definitionsverweis {offen}{}{,}
eine
\definitionsverweis {differenzierbare Hyperfläche}{}{}
und sei
\maabb {\gamma} {I =[a,b]} {Y
} {}
eine
\definitionsverweis {geodätische Kurve}{}{.}
Zeige, dass der
\definitionsverweis {Paralleltransport}{}{}
längs $\gamma$ den tangentialen Vektor

\mavergleichskette
{\vergleichskette
{ \gamma'(a)
}
{ \in }{ T_{\gamma(a)}Y
}
{  }{
}
{    }{
}
{   }{
}
}
{}{}{}
in den tangentialen Vektor

\mavergleichskette
{\vergleichskette
{ \gamma'(b)
}
{ \in }{ T_{\gamma(b)}Y
}
{  }{
}
{    }{
}
{   }{
}
}
{}{}{}
überführt.

}
{} {}

\inputaufgabe
{}
{

Es sei $Y$ die zweidimensionale
\definitionsverweis {Einheitssphäre}{}{}
und
\maabb {\gamma} {[0,2\pi]} {Y
} {}
die
\definitionsverweis {Bogenparametrisierung}{}{}
eines
\definitionsverweis {Großkreises}{}{}
mit

\mavergleichskette
{\vergleichskette
{ \gamma(0)
}
{ = }{ \gamma(2 \pi)
}
{ = }{ P
}
{ \in   }{ Y
}
{   }{
}
}
{}{}{.}
Zeige, dass der zugehörige
\definitionsverweis {Paralleltransport}{}{}
die Identität auf $T_PY$ ist.

}
{} {}

\inputaufgabe
{}
{

Es sei

\mavergleichskettedisp
{\vergleichskette
{ Y
}
{ =} { \left\{  \begin{pmatrix}  x  \\ y \\ z   \end{pmatrix}   \mid   x^2+y^2+z^2 = 1        \right\}
}
{ } {
}
{ } {
}
{ } {
}
}
{}{}{}
die zweidimensionale
\definitionsverweis {Einheitssphäre}{}{,}

\mavergleichskette
{\vergleichskette
{ P
}
{ = }{ \begin{pmatrix}  { \frac{   1  }{  2 } }  \\ 0 \\  { \frac{   \sqrt{3}  }{  2 } }   \end{pmatrix}
}
{  }{
}
{    }{
}
{   }{
}
}
{}{}{}
und
\maabbele {\gamma} {[0,2\pi]} {Y
} { t } { \begin{pmatrix}  { \frac{   1  }{  2  } } \cos  t  \\ { \frac{   1  }{  2  } }  \sin  t \\  { \frac{   \sqrt{3}  }{  2  } }   \end{pmatrix}
} {.}
Beschreibe den
\definitionsverweis {Paralleltransport}{}{}
\maabb {\Psi_\gamma} { T_PY} {T_PY
} {}
bezüglich einer geeigneten Basis.

}
{} {}

\inputaufgabe
{}
{

Es sei

\mavergleichskette
{\vergleichskette
{ Y
}
{ \subseteq }{ W
}
{ \subseteq }{ \R^n
}
{    }{
}
{   }{
}
}
{}{}{,}
$W$
\definitionsverweis {offen}{}{,}
eine
\definitionsverweis {differenzierbare Hyperfläche}{}{}
und seien

\mavergleichskette
{\vergleichskette
{ P,Q
}
{ \in }{ Y
}
{  }{
}
{    }{
}
{   }{
}
}
{}{}{}
Punkte mit den
\definitionsverweis {Tangentialräumen}{}{}
\mathkor {} {T_PY} {und} {T_QY} {.}
Wir betrachten die Abbildung
\maabbeledisp {\varphi_{P,Q}} {T_PY} { T_QY
} { v } { \pi_Q(v)
} {,}
die einen Tangentialvektor in $P$ als einen Vektor im umgebenden Raum $\R^n$ auffasst und anschließend auf $T_QY$ orthogonal projiziert.
\aufzaehlungvier{Ist
\mathl{\varphi_{P,Q}}{} linear?
}{Ist
\mathl{\varphi_{P,Q}}{} bijektiv?
}{Ist
\mathl{\varphi_{P,Q}}{} eine Isometrie?
}{Gibt es eine Beziehung zwischen
\mathl{\varphi_{P,Q}}{} und dem
\definitionsverweis {Paralleltransport}{}{} $\Psi_\gamma$ zu einer differenzierbaren Kurve
\maabb {\gamma} {[a,b]} { Y
} {}
mit

\mavergleichskette
{\vergleichskette
{ \gamma(a)
}
{ = }{ P
}
{  }{
}
{    }{
}
{   }{
}
}
{}{}{}
und

\mavergleichskette
{\vergleichskette
{ \gamma(b)
}
{ = }{ Q
}
{  }{
}
{    }{
}
{   }{
}
}
{}{}{.}
}

}
{} {}

Es sei

\mavergleichskette
{\vergleichskette
{ W
}
{ \subseteq }{ \R^n
}
{  }{
}
{    }{
}
{   }{
}
}
{}{}{}
\definitionsverweis {offen}{}{,}
\maabb {h} { W } { \R
} {}
eine
\definitionsverweis {stetig differenzierbare Funktion}{}{}
und

\mavergleichskette
{\vergleichskette
{ Y
}
{ = }{ h^{-1}(c)
}
{  }{
}
{    }{
}
{   }{
}
}
{}{}{}
die
\definitionsverweis {Faser}{}{}
zu

\mavergleichskette
{\vergleichskette
{ c
}
{ \in }{ \R
}
{  }{
}
{    }{
}
{   }{
}
}
{}{}{,}
wobei $h$ in jedem Punkt von $Y$
\definitionsverweis {regulär}{}{}
sei. Zu einem Punkt

\mavergleichskette
{\vergleichskette
{ P
}
{ \in }{ Y
}
{  }{
}
{    }{
}
{   }{
}
}
{}{}{}
nennt man die
\definitionsverweis {Untergruppe}{}{}

\mavergleichskette
{\vergleichskette
{ H
}
{ \subseteq }{ \operatorname{Iso} \left(  T_PY  \right)
}
{  }{
}
{    }{
}
{   }{
}
}
{}{}{,}
die aus allen durch stückweise differenzierbaren
\definitionsverweis {geschlossenen Wegen}{}{}
$\gamma$ auf $Y$ von $P$ nach $P$ induzierten
\definitionsverweis {Paralleltransporten}{}{}
\maabbdisp {\Psi_\gamma} {T_PY} {T_PY
} {}
besteht, die
\definitionswort {Holonomiegruppe}{}
zu $P$.

\inputaufgabe
{}
{

Es sei

\mavergleichskette
{\vergleichskette
{ W
}
{ \subseteq }{ \R^n
}
{  }{
}
{    }{
}
{   }{
}
}
{}{}{}
\definitionsverweis {offen}{}{,}
\maabb {h} { W } { \R
} {}
eine
\definitionsverweis {stetig differenzierbare Funktion}{}{}
und

\mavergleichskette
{\vergleichskette
{ Y
}
{ = }{ h^{-1}(c)
}
{  }{
}
{    }{
}
{   }{
}
}
{}{}{}
die
\definitionsverweis {Faser}{}{}
zu

\mavergleichskette
{\vergleichskette
{ c
}
{ \in }{ \R
}
{  }{
}
{    }{
}
{   }{
}
}
{}{}{,}
wobei $h$ in jedem Punkt von $Y$
\definitionsverweis {regulär}{}{}
sei. Zeige, dass die
\definitionsverweis {Holonomiegruppe}{}{}
zu einem Punkt

\mavergleichskette
{\vergleichskette
{ P
}
{ \in }{ Y
}
{  }{
}
{    }{
}
{   }{
}
}
{}{}{}
in der Tat eine
\definitionsverweis {Untergruppe}{}{}
der
\definitionsverweis {Isometriegruppe}{}{}
ist.

}
{} {}

\inputaufgabe
{}
{

Es sei

\mavergleichskette
{\vergleichskette
{ Y
}
{ \subseteq }{ \R^n
}
{  }{
}
{    }{
}
{   }{
}
}
{}{}{}
eine
\definitionsverweis {Hyperebene}{}{.}
Bestimme die
\definitionsverweis {Holonomiegruppe}{}{}
zu

\mavergleichskette
{\vergleichskette
{ P
}
{ \in }{ Y
}
{  }{
}
{    }{
}
{   }{
}
}
{}{}{.}

}
{} {}

  \zwischenueberschrift{Aufgaben zum Abgeben}

\inputaufgabe
{2}
{

Wir betrachten auf der zweidimensionalen
\definitionsverweis {Einheitssphäre}{}{}
$Y$ das
\definitionsverweis {tangentiale Vektorfeld}{}{}

\mavergleichskettedisp
{\vergleichskette
{ F \begin{pmatrix}  x  \\ y \\ z   \end{pmatrix}
}
{ =} { \begin{pmatrix}  y  \\ -x\\  0   \end{pmatrix}
}
{ } {
}
{ } {
}
{ } {
}
}
{}{}{.}
Bestimme für den Punkt

\mavergleichskette
{\vergleichskette
{ R
}
{ = }{ \begin{pmatrix}  { \frac{   1  }{  2 } }  \\ { \frac{   1  }{  2 } } \\  { \frac{   1  }{  \sqrt{2}  } }   \end{pmatrix}
}
{ \in }{ Y
}
{    }{
}
{   }{
}
}
{}{}{}
eine Matrix für die Abbildung
\maabbeledisp {} {T_RY} { T_RY
} { v } { \nabla_v F
} {,}
bezüglich einer geeigneten
\definitionsverweis {Basis}{}{.}

}
{} {}

\inputaufgabe
{4}
{

Es sei

\mavergleichskettedisp
{\vergleichskette
{ h(x,y,z)
}
{ =} { x^2+y^2-2z^7
}
{ } {
}
{ } {
}
{ } {
}
}
{}{}{}
und

\mavergleichskette
{\vergleichskette
{ Y
}
{ = }{ h^{-1}(3)
}
{  }{
}
{    }{
}
{   }{
}
}
{}{}{.}
Es sei
\maabbeledisp {\gamma} { [1,5] } { Y
} { t } { \begin{pmatrix}  \sqrt{3  -4t^2+2t^7  } \\ 2t\\ t   \end{pmatrix}
} {,}
ein differenzierbarer Weg auf $Y$. Bestimme die gewöhnliche Differentialgleichung in der Form von
[[Hyperfläche/Kurve/Anfangstangentenvektor/Paralleles Vektorfeld/Differentialgleichung/Explizite Gestalt/Bemerkung|Bemerkung 6.10]]
für die
\definitionsverweis {parallelen Vektorfelder}{}{}
längs $\gamma$.

}
{} {}

\inputaufgabe
{4}
{

Zeige, dass die
\definitionsverweis {Holonomiegruppe}{}{}
auf der
\definitionsverweis {Einheitssphäre}{}{}
gleich der
\definitionsverweis {Isometriegruppe}{}{}
ist.

}
{} {}

\inputaufgabe
{4}
{

Es sei

\mavergleichskette
{\vergleichskette
{ W
}
{ \subseteq }{ \R^n
}
{  }{
}
{    }{
}
{   }{
}
}
{}{}{}
\definitionsverweis {offen}{}{,}
\maabb {h} { W } { \R
} {}
eine
\definitionsverweis {stetig differenzierbare Funktion}{}{}
und

\mavergleichskette
{\vergleichskette
{ Y
}
{ = }{ h^{-1}(c)
}
{  }{
}
{    }{
}
{   }{
}
}
{}{}{}
die
\definitionsverweis {Faser}{}{}
zu

\mavergleichskette
{\vergleichskette
{ c
}
{ \in }{ \R
}
{  }{
}
{    }{
}
{   }{
}
}
{}{}{,}
wobei $h$ in jedem Punkt von $Y$
\definitionsverweis {regulär}{}{}
sei. Es sei
\maabbdisp {L} {\R^n } { \R^n
} {}
eine
\definitionsverweis {lineare Isometrie}{}{,}

\mavergleichskette
{\vergleichskette
{ W'
}
{ = }{ L^{-1}(W)
}
{  }{
}
{    }{
}
{   }{
}
}
{}{}{}
und

\mavergleichskettedisp
{\vergleichskette
{ Z
}
{ =} { L^{-1}(Y)
}
{ =} { (h \circ L)^{-1}(c)
}
{ } {
}
{ } {
}
}
{}{}{.}
Man erläutere, wie sich die Konzepte
\definitionsverweis {paralleles Vektorfeld}{}{,}
\definitionsverweis {Paralleltransport}{}{}
und
\definitionsverweis {Holonomiegruppe}{}{}
unter $L$ verhalten.

}
{} {}

 [[Kategorie:Latexseite]]
